\section{The Weil actions in the doubling calculation}

% BEGIN W87N1
\paragraph{W87N1. The scalar on the auxiliary hyperbolic target.}
Let $H$ be the orthogonal sum of the $t$ hyperbolic planes in
Y87R1, with splitting character $\chi_H=1$. For the given oriented
Weil datum on $(A,H)$, its restriction to $U(A)(\mathbb A_F)$ in
the polarization by a rational maximal isotropic half of $H$ acts
by change of variables on $t$ copies of the dual of $A$.
In the polarization by one factor of a finite split place it acts
by change of variables on the $2t$ rows, with the positive square
root of the change-of-measure factor. Both statements concern the
actual action on the theta space.

To see how the scalar is determined, take two actions on that space
which implement the same Heisenberg automorphisms and preserve the
Haar norm. For a fixed group element $g$ the inverse of the second
operator times the first commutes with all Heisenberg translations
and preserves the norm. CW15 makes it a scalar $c(g)$ of modulus one.
The theta space contains a nonzero vector (for instance the Gaussian
times a finite lattice test in a polarization). Since each operator is
invertible, its scalar is unique and the group law gives
$c(gg')=c(g)c(g')$. Translate a nonzero theta vector by a Heisenberg
element to get $f$ with $f(0)\ne0$. At any $g_0$ choose $w$ with
$(\omega'(g_0)f)(w)\ne0$. In a neighborhood of $g_0$ the scalar is
\[
          c(g)=\frac{(\omega(g)f)(w)}{(\omega'(g)f)(w)} .
\]
Both evaluations are continuous by the pointwise continuity clause
of the Weil datum, or by the explicit Schwartz action for the
comparison operator. Thus $c$ is continuous. Rational automorphy
at zero for both actions, applied to $f$, gives $c(\gamma)=1$ for
$\gamma\in U(A)(F)$.

At a finite split place $v$ choose the summand $e$ so that the
conjugate summand acts by modulations in the local Schr\"odinger
model. Evaluation there at zero followed by any nonzero evaluation
on the other tensor factors gives a linear functional invariant
under these modulations, nonzero on a theta test. It may be extended
linearly to the space of all functions as required in the native
axiom, since that axiom tests invariance only on theta vectors.
For $g$ supported at $v$ the native split axiom acts on this
functional by a positive real times $\chi_H(x)^{-1}$, with $x$
the determinant on the $e$-component extended by one. More
generally the phase for a pair $(g,h)$ is
$\chi_H(x)^{-1}\chi_A(y)$ for the analogous determinant $y$ of $h$.
If a comparison action has this same phase, its ratio with the
given action at $g$ is positive real and has modulus one.
Consequently $c(g)=1$ at every finite split place.

Here is the character argument extending that conclusion to the
adeles. For $d=\dim_E A\ge2$ the special unitary group is simply
connected and absolutely almost simple. At a split finite place
$v_0$ it is $\mathrm{SL}_d(F_{v_0})$, which is noncompact. Strong
approximation away from $v_0$, rational triviality and triviality
at $v_0$ kill $c$ on the entire adelic special unitary group.
For $d=1$ the determinant already identifies the unitary group
with the norm-one torus. For every positive $d$, scaling a
nonisotropic rational line and fixing its orthogonal complement
gives a rational section of the determinant onto that torus.
Pull the character there back by $x\mapsto x/\bar x$. This gives
a unitary Hecke character $\xi$ of $E$, trivial at the places
above split places of $F$. The map is surjective on adelic
points by local Hilbert 90: at almost every inert place a preimage
can be made integral of valuation zero by multiplying by an
$F_v$-uniformizer to an integral power, and the split case has
the section $z\mapsto(z,1)$. Omit the finitely many ramified
places. The quotient $L_E^S(s,\xi)/\zeta_E^S(s)$ has only inert
Euler factors, whose prime norms are $q_v^2$. Its logarithmic
series converges absolutely near $s=1$. The quotient is nonzero
there, so $L_E^S(s,\xi)$ has a pole at $1$. The pole theorem for
Hecke $L$-functions gives $\xi=1$. Surjectivity now gives $c=1$.
The same argument applied separately to each factor gives uniqueness
on the whole pair whenever both factors have the same split phases.

For the asserted auxiliary action, take the explicit pullback of
Schwartz functions in the polarization from $H$. The isometry of
$A$ preserves its $F$-volume, as its $E$-determinant has norm one.
Thus this pullback preserves the norm and implements the
transformation of translations and modulations. Periodization sends
it to an action on the oriented theta space. Concretely, if $X$ is
the chosen rational configuration space and $\rho$ is the Weyl
action for the native alternating form and character, this transform
sends $\phi$ to
\[
                  w\longmapsto\sum_{a\in X(F)}(\rho(w)\phi)(a).
\]
Weyl multiplication gives the native Heisenberg translation, with
factor $\psi(\langle w,w_0\rangle/2)$, and the inverse takes the
Fourier coefficients along the complementary rational Lagrangian
modulo rational points. The characters of this compact quotient
are indexed by the rational configuration space. Quasiperiodicity
identifies their coefficients with rational translates of a single
function. Finite-level invariance makes the result locally constant
with bounded finite support (translation invariance gives the level,
modulation invariance the support). Repeated differentiation on the
compact real torus gives rapid decay of these coefficients, also
after differentiation in the real configuration coordinates. Thus
the inverse is Schwartz. Fourier inversion and Parseval give the
inverse identities and the Haar norm, up to the same positive
measure constant on both sides. Rational elements of $U(A)$
permute the summands at $w=0$, proving rational automorphy.
The periodized evaluations are continuous by Schwartz convergence,
locally uniform under changes of variables.

Finally, changing at $v$ from the hyperbolic polarization of $H$
to the split $e$-polarization exchanges the variables of one
isotropic half with its dual by partial Fourier transform.
The Fourier transform of a linear change of variables is the dual
linear change with the positive absolute Jacobian. The original
pullback had scalar one, so evaluation at zero in the new model
has positive phase for $g$ at $v$, exactly $\chi_H(x)^{-1}=1$.
The character argument identifies these actions for every $g$.
In particular the integral lattice test away from the exceptional
places and the real Gaussian in the positive metric of $A$ have
the required source invariance.
% END W87N1

% BEGIN W87N2
\paragraph{W87N2. The section and the target kernel in the same standard chart.}
Write $U=A\oplus -A$, $n=\dim_E A$, and let $T$ be either
$W$ or $B=W\perp H$, with $b=\dim_E T$. The doubled datum has
characters $(1,\chi_W)$ for both values of $T$. The source
characters cancel on doubling. On adjoining $H$ the target
character is $\chi_W\chi_H=\chi_W$. Indeed on the tensor
product for an orthogonal sum Heisenberg translations factor, the
Haar norms multiply, and rational evaluation factors on products.
The split phases on $U(A)$ multiply in the same way. Applying the
scalar argument of W87N1 to this restriction on the tensor product
identifies the action there. Evaluation on tensor products suffices:
finite-place functions are finite sums of rectangle functions and
the products of real Hermite functions are complete in Schwartz
seminorms.

Choose the diagonal polarization of $U$. The Schr\"odinger map on
a doubled theta function $f$ takes the zero Fourier coefficient
along the diagonal Lagrangian after translation by the configuration
variable $x\in\operatorname{Mat}_{b,n}(\mathbb A_E)$. Call the
result $\phi(x)$. Fourier inversion reconstructs the theta vector
by the periodization of W87N1. With covolume one, $\phi(0)$ is
exactly the diagonal integral defining the native Siegel--Weil
section. In particular the transport applies to every Schwartz
test occurring in a residue or in an Ikeda map. It has an exact
formula also on the target group: for all $g\in U(A\oplus-A)(\mathbb A_F)$
and $h\in U(T)(\mathbb A_F)$ the original theta kernel equals
\[
       \sum_{a\in\operatorname{Mat}_{b,n}(E)}
                     (\omega_{\rm Sch}(g)\phi)(h^{-1}a).
\]
Indeed $h$ acts on the configuration variable by inverse
left multiplication, whose modulus is one; the source splitting
character here is one. The scalar follows either from W87N1 or
from split evaluation and rational triviality applied to $U(T)$.
Periodize after this action to obtain the formula.

For the source group use a standard unitary chart with nonzero skew
scalar $\mu$ and put $\lambda=-\xi_0/\mu\in F^\times$,
$\psi'(z)=\psi(\lambda z)$. Let
$[x,y]=x^*G_Ty$, and use $\operatorname{Tr}_{E/F}$ on scalar
matrix traces. The three transported operators in this chart are
\begin{align*}
 m(a):\quad&\phi(x)\longmapsto
       \chi_W(\det a)^{-1}|\det a|_E^{b/2}\phi(xa),\\
 n(u):\quad&\phi(x)\longmapsto
       \psi'\bigl(\tfrac12\operatorname{Tr}_{E/F}
                   \operatorname{tr}(u[x,x])\bigr)\phi(x),\\
 w:\quad&\phi(x)\longmapsto
       \int\phi(y)\psi'\bigl(\operatorname{Tr}_{E/F}
                   \operatorname{tr}([x,y])\bigr)\,dy ,
       \qquad w=\begin{pmatrix}0&I\\-I&0\end{pmatrix}.
\end{align*}
For the Levi, translations determine the change in $x$ and Haar
isometry gives the half modulus. The remaining unitary character
factors through the determinant: elementary unipotents generating
the local special linear group are commutators in the general
linear group, and local factors generate a dense subgroup. At
split places the zero-coefficient functional gives the inverse
target character on $\det a$. Rational automorphy and the Hecke
character argument in W87N1 extend this phase globally.
For the unipotent, covariance gives the character of the quadratic
Gram form; a possible extra scalar is an additive character of the
unipotent group invariant under Levi conjugation. Conjugating by
the rational central scalar $2$ sends $u$ to $4u$, so the character
is trivial on every $3u$ and hence is trivial. Substituting
$-\xi_0/\mu$ into the native half trace gives exactly the displayed
sign. For $w$, Fourier inversion exchanges translations and
modulations, fixing the operator up to a scalar. The chart's $w$
is rational; Poisson summation on the rational matrices, with the
self-dual covolume-one measure, gives equality of the theta
functional before and after the displayed Fourier operator. Its
value on a Gaussian times a suitable finite indicator is nonzero,
so rational automorphy fixes that scalar as one. These equations
also identify the operators at arbitrary adelic elements:
conjugating each local Siegel unipotent by this $w$ gives its
opposite at that place. The two parabolics generate the local
quasi-split unitary group by Bruhat decomposition; their local
subgroups and the archimedean factors generate a dense subgroup
of the adeles. Pointwise continuity gives equality on its closure.
Thus both complementary targets use the standard action with the
same $\psi'$ and the same inducing character $\chi_W^{-1}$.
At the section point $s=(b-n)/2$ the factor
$|\det a|_E^{b/2}$ is $|\det a|_E^{n/2+s}$, including the
half modulus of the Siegel parabolic.

The sign in the two-copy pullback can be checked independently.
Negating the form reverses the Heisenberg character. Accordingly
the second test is $x\mapsto\overline{f(-x)}$. At a split place the
conjugated source phase on that copy is $\chi_W(x)$, whereas
the doubled action prescribes $\chi_W(x)^{-1}$. Here
$x/\bar x=\det g_2$. The rank restriction of $\chi_W$ implies
$\chi_W(x\bar x)=1$ because the quadratic character kills norms.
Thus the ratio is $\chi_W(x)^{-2}=\chi_W(\det g_2)^{-1}$.
The two actions after this twist have equal rational and split
normalizations. By W87N1 they coincide globally on this factor.
The doubled kernel therefore acquires
$\chi_W(\det g_2)^{-1}$, cancelled by the
$\chi_W(\det g_2)$ in the pullback of Y87R4.
% END W87N2
